\subsection{取值于 Banach 空间中的函数的叠积分}

定理 1. --- 设 f 是取值于 Banach 空间 F 或 \(\overline{R}\) 中的一个函数，并设 H 是由 \(t \in T\) （使 f 非 \(\lambda_{t}\) -可积者）组成的集。

\begin{enumerate}
\def\labelenumi{\alph{enumi})}
\tightlist
\item
  若 \(\mathbf{f}\) 是 \(\nu\) -可积的，则 \(\mathrm{H}\) 是局部 \(\mu\) -可忽略的，函数 \(t \mapsto \int \mathbf{f}(x)d\lambda_t(x)\) （对 \(t \notin \mathrm{H}\) 有定义）是本质 \(\mu\) -可积的，且
\end{enumerate}

\[
\int \mathbf {f} (x) d \nu (x) = \int d \mu (t) \int \mathbf {f} (x) d \lambda_ {t} (x).\tag{11}
\]

\begin{enumerate}
\def\labelenumi{\alph{enumi})}
\setcounter{enumi}{1}
\item
  若 \(\mathbf{f}\) 是 \(\nu\) -可积的且 \(\Lambda\) 是淡连续的，则此外 \(\mathbf{H}\) 还是 \(\mu\) -可忽略的，且函数 \(t \mapsto \int \mathbf{f}(x) d\lambda_t(x)\) （对 \(t \notin \mathbf{H}\) 有定义）是 \(\mu\) -可积的。
\item
  若 \(\lambda_t^\bullet(1) < +\infty\) 局部 \(\mu\) -几乎处处成立，则对本质 \(\nu\) -可积函数 f，a) 的结论仍成立。
\end{enumerate}

我们先来建立 a)（相应地，b)）。当 f 是正数值函数时，这一命题成立（命题 5）；若 f 是取值于 \(\overline{R}\) 的可积函数，则这一结果可以应用于正函数 \(f^{+}\) 与 \(f^{-}\) ，从而由减法立即推广到 f。剩下要处理取值于 F 的函数的情形。设 H 是 \(\mathcal{L}_{\mathrm{F}}^{1}(\nu)\) 中由 \(\mathcal{H}(\mathrm{X})\) 里的函数以 F 中元素为系数的线性组合所构成的子空间；关于实函数的结果立即蕴涵该命题对 H 的元素成立。而 H 在 \(\mathcal{L}_{\mathrm{F}}^{1}(\nu)\) 中稠密；因此，对每个 \(\mathbf{f} \in \mathcal{L}_{\mathrm{F}}^{1}(\nu)\) ，存在由 H 中元素组成的一个序列 \((\mathbf{f}_{n})\) ，它具有下列性质：

\begin{enumerate}
\def\labelenumi{\arabic{enumi})}
\item
  序列 \((\mathbf{f}_n)\) 平均收敛于 \(\mathbf{f}\) （在 \(\mathcal{L}_{\mathrm{F}}^{1}(\nu)\) 中），且 \(\nu\) -几乎处处收敛；
\item
  函数 \(g = |\mathbf{f}_0| + \sum_{n\in \mathbf{N}}|\mathbf{f}_{n + 1} - \mathbf{f}_n|\) 满足 \(\nu^{*}(g) < +\infty\) （Ch. IV, §3, No.~4, 定理 3）。
\end{enumerate}

设 \(N_{1}\) 是由 \(t \in T\) （使 \(\lambda_{t}^{*}(g) = +\infty\) 者）组成的集；由公式 (6)（相应地，(7)），\(N_{1}\) 是局部 \(\mu\) -可忽略的（相应地，\(\mu\) -可忽略的）。对 \(t \notin N_{1}\) ，诸 \(f_{n}\) 属于 \(\mathcal{L}_{F}^{1}(\lambda_{t})\) ，序列 \((f_{n})\) \(\lambda_{t}\) -几乎处处收敛，且对 \(\mathcal{L}_{F}^{1}(\lambda_{t})\) 中的平均收敛拓扑也收敛（Ch. IV, §3, No.~3, 命题 6）。设 M 是由 \(x \in X\) （使 \(f_{n}(x)\) 不收敛于 \(f(x)\) 者）组成的集：由于 M 是 \(\nu\) -可忽略的，集 \(N_{2}\) （由 \(t \in T\) （使 M 非 \(\lambda_{t}\) -可忽略者）组成）由命题 3 的推论 1 是局部 \(\mu\) -可忽略的（相应地，\(\mu\) -可忽略的）。

设 \(t\) 不属于 \(\mathrm{N}_1 \cup \mathrm{N}_2\) ；序列 \((\mathbf{f}_n)\) 在 \(\mathcal{L}_{\mathrm{F}}^{1}(\lambda_t)\) 中平均收敛，且 \(\lambda_t\) -几乎处处收敛于 \(\mathbf{f}\) 。因此 \(\mathbf{f} \in \mathcal{L}_{\mathrm{F}}^{1}(\lambda_t)\) 且 \(\int \mathbf{f} d\lambda_t = \lim_{n \to \infty} \int \mathbf{f}_n d\lambda_t\) （Ch. IV, §4, No.~1）。于是命题中的集 H 含于 \(N_{1} \cup N_{2}\) ；从而它是局部 \(\mu\) -可忽略的（相应地，\(\mu\) -可忽略的）。另一方面，函数 \(t \mapsto \int f d\lambda_{t}\) 局部 \(\mu\) -几乎处处等于一列 \(\mu\) -可测函数的极限；因此它是 \(\mu\) -可测的。最后，对每个 \(t \notin N_{1} \cup N_{2}\) 及每个 n，有

\[
\left| \int \mathbf {f} _ {n} (x) d \lambda_ {t} (x) \right| \leqslant \int^ {*} g (x) d \lambda_ {t} (x)
\]

这依据不等式 \(|\mathbf{f}_n| \leqslant g\) 及 Ch. IV, §4, No.~2 的命题 2。而函数 \(t \mapsto \int^{*} g(x) d\lambda_t(x)\) 由命题 5 是本质 \(\mu\) -可积的（相应地，\(\mu\) -可积的）。因此可以应用 Lebesgue 定理，它给出

\[
\int d \mu (t) \int \mathbf {f} (x) d \lambda_ {t} (x) = \lim _ {n \rightarrow \infty} \int d \mu (t) \int \mathbf {f} _ {n} (x) d \lambda_ {t} (x) = \lim _ {n \rightarrow \infty} \int \mathbf {f} _ {n} (x) d \nu (x).
\]

由于 \(\int \mathbf{f}_n(x) d\nu(x)\) 趋于 \(\int \mathbf{f}(x) d\nu(x)\) （当 \(n\) 趋于 \(\infty\) 时；依据对序列 \((\mathbf{f}_n)\) 所作的假设），关系式 (11) 得出，从而我们证明了 a)（相应地，b)）。

现在假设 \(\lambda_t^\bullet(1) < +\infty\) 局部 \(\mu\) -几乎处处成立，且 \(\mathbf{g}\) 是本质 \(\nu\) -可积函数。设 \(\mathbf{f}\) 是一个 \(\nu\) -可积函数，使得 \(\mathbf{g} = \mathbf{f}\) 局部 \(\nu\) -几乎处处成立（§1, No.~3）。于是 \(\mathbf{g} = \mathbf{f}\) 对 \(\lambda_t\) 几乎处处成立，除去由 \(t\) 组成的一个局部 \(\mu\) -可忽略集 \(\mathrm{P}\) 的情形（命题 3 的推论 1 c)）。因此 \(\int \mathbf{g} d\lambda_t = \int \mathbf{f} d\lambda_t\) 对所有 \(t \notin \mathrm{P} \cup \mathrm{H}\) 成立，这就完成了证明。

注. --- 设 \(\Lambda: t \mapsto \lambda_{t}\) 是一个 \(\mu\) -适宜映射（T 到 \(\mathcal{M}_{+}(X)\) 中的）。如果一个映射 \(\Lambda': t \mapsto \lambda_{t}'\) （T 到 \(\mathcal{M}_{+}(X)\) 中的）与 \(\Lambda\) 局部 \(\mu\) -几乎处处相等，则由定义立即可见 \(\Lambda'\) 也是 \(\mu\) -适宜的，且 \(\Lambda\) 与 \(\Lambda'\) 有相同的积分。现在设 \(H: t \mapsto \eta_{t}\) 是一个取值于 \(\mathcal{M}_{+}(X)\) 中的函数，它只在局部 \(\mu\) -几乎处处有定义；我们仍称 H 是 \(\mu\) -适宜的，如果它局部 \(\mu\) -几乎处处等于一个映射 \(\Lambda: t \mapsto \lambda_{t}\) ，后者处处有定义且 \(\mu\) -适宜。这时我们置 \(\int \eta_{t} d\mu(t) = \int \lambda_{t} d\mu(t)\) ，这一定义不依赖于所用的函数 \(\Lambda\) 。前面各 No. 中证明的命题均可推广到局部 \(\mu\) -几乎处处有定义的 \(\mu\) -适宜函数，验证工作留给读者。

